\begin{abstract}
RabbitMQ is a widely used open-source message broker and its quorum queues use the Raft consensus algorithm to replicate messages across cluster nodes for data safety and high availability. Horizontal scaling is a common strategy to increase message throughput, but it is not obvious whether adding nodes delivers proportional gains when every message must be replicated. Additionally, Raft requires a majority of nodes to acknowledge every write before it is committed, creating a trade-off between fault tolerance and latency that makes it unclear how latency behaves as the cluster grows. We investigate this trade-off by developing a Go CLI benchmark tool and running three experiments, each repeated three times, across cluster sizes of 3, 6 and 9 nodes on Azure infrastructure, measuring maximum throughput with fixed quorum size and end-to-end latency with quorum size equal to the cluster size. Our results show that throughput scales near-linearly (3.11$\times$ at 9 nodes) when quorum size is fixed at 3, while end-to-end mean latency increases from 3.4\,ms to 4.3\,ms and P99 latency rises from 4.6\,ms to 6.0\,ms with growing inter-run variability as more nodes participate in Raft replication. Under heavy load, latency degrades by up to 146$\times$ compared to the ideal baseline. Our results suggest that adding nodes with a fixed quorum size is effective for scaling throughput, but increasing the quorum size carries a significant latency cost that should be weighed carefully against the additional fault tolerance it provides.
\end{abstract}
